\documentclass[12pt]{article}
\usepackage[T1]{fontenc}
\usepackage[utf8]{inputenc}
\usepackage[brazil]{babel}
\usepackage[margin=1in]{geometry}
\setlength{\parindent}{0pt}
\begin{document}
\section*{Tema 25}
Processo: PEDILEF 2005.82.00.505195-9/ PB\\
Situacao: Julgado\\
Ramo: DIREITO PREVIDENCIÁRIO\\
Questao: Saber qual a tábua de mortalidade elaborada pelo IBGE a ser utilizada no cálculo do fator previdenciário.\\
Tese: Para o cálculo do fator previdenciário deve ser observada a tábua de mortalidade elaborada pelo IBGE vigente na data do requerimento do benefício previdenciário, e não aquela utilizada anteriormente, quando preenchidos os requisitos necessários à concessão da aposentadoria.\\
Fonte: https://www.cjf.jus.br/phpdoc/virtus/
\end{document}
